\documentclass[a4paper]{arrowhead}

\usepackage[yyyymmdd]{datetime}
\usepackage{etoolbox}
\usepackage[utf8]{inputenc}
\usepackage[T1]{fontenc}
\usepackage{multirow}
\usepackage{longtable}
\usepackage{tikz}
\usepackage{adjustbox}
\AtBeginEnvironment{longtable}{\small}
\BeforeBeginEnvironment{tikzpicture}{\begin{adjustbox}{max width=\textwidth}}
\AfterEndEnvironment{tikzpicture}{\end{adjustbox}}
\usetikzlibrary{positioning,shapes.multipart,arrows.meta}

\renewcommand{\dateseparator}{-}
\setlength{\emergencystretch}{3em}
\lstset{language=HTTP,breakatwhitespace=false,columns=fullflexible}

%% Special references
\newcommand{\mref}[2][]{{\textcolor{ArrowheadPurple}{\hyperref[sec:model:#2]{\ifstrempty{#1}{#2}{#1}}}}}
\newcommand{\oref}[1]{{\textcolor{ArrowheadBlue}{\hyperref[sec:operations:#1]{#1}}}}
\newcommand{\pdef}[1]{{\textcolor{ArrowheadGrey}{#1\label{sec:model:primitives:#1}\label{sec:model:primitives:#1s}\label{sec:model:primitives:#1es}}}}
\newcommand{\pref}[1]{{\textcolor{ArrowheadGrey}{\hyperref[sec:model:primitives:#1]{#1}}}}

% as template_sd.tex; the input and output arguments are passed already formatted (\mref or \pref)
\newrobustcmd\osubsection[3]{
  \addtocounter{subsection}{1}
  \addcontentsline{toc}{subsection}{\protect\numberline{\thesubsection}operation \textcolor{ArrowheadBlue}{#1}}
  \subsection*{
    \thesubsection\quad
    operation
    \textcolor{ArrowheadBlue}{#1}
    (\notblank{#2}{#2}{})
    \notblank{#3}{: #3}{}
  }
  \label{sec:operations:#1}
}
\newrobustcmd\msubsection[2]{
  \addtocounter{subsection}{1}
  \addcontentsline{toc}{subsection}{\protect\numberline{\thesubsection}#1 \textcolor{ArrowheadPurple}{#2}}
  \subsection*{\thesubsection\quad#1 \textcolor{ArrowheadPurple}{#2}}
  \label{sec:model:#2} \label{sec:model:#2es}
}
%%

\usepackage{xurl}
\begin{document}
\sloppy

%% Arrowhead Document Properties
\ArrowheadTitle{serviceOrchestrationPlantManagement}
\ArrowheadServiceID{serviceOrchestrationPlantManagement}
\ArrowheadType{Service Description}
\ArrowheadTypeShort{SD}
\ArrowheadVersion{5.3.0}
\ArrowheadDate{2026-11-09}
\ArrowheadAuthor{Jerker Delsing}
\ArrowheadStatus{REVIEW DRAFT}
\ArrowheadContact{jerker.delsing@ltu.se}
\ArrowheadFooter{\href{www.arrowhead.eu}{www.arrowhead.eu}}
\ArrowheadSetup
%%

%% Front Page
\begin{center}
  \vspace*{1cm}
  \huge{\arrowtitle}

  \vspace*{0.2cm}
  \LARGE{\arrowtype}
  \vspace*{1cm}

  \Large{Service ID: \textit{"\arrowid"}}

  \vspace*{\fill}

  % Front Page Image

  \vspace*{1cm}
  \vspace*{\fill}

  % Front Page Abstract
  \begin{abstract}
    This document describes the \texttt{service\allowbreak{}Orchestration\allowbreak{}Plant\allowbreak{}Management} service, version 1.0.0, produced by the DynamicServiceOrchestration core system of Arrowhead (documentation edition 5.3.0).
  \end{abstract}

  \vspace*{1cm}
\end{center}

\newpage
%%

%% Table of Contents
\tableofcontents
\newpage
%%

\section{Documentation edition and implementation scope}
This is the 5.3.0 documentation edition, dated 2026-11-09. It carries forward the system and service contracts from the retained 5.2.2 documentation and identifies the current primary implementation sources. The former \path|documentation/v5.2.2| folder is retained unchanged.

Document version 5.3.0 is distinct from executable, dependency, registered system/service and wire-contract versions. Values still defined by the source (including 5.2.2 artefact versions, 1.0.0 service versions and \path|SysMonitor-AMQP-IDD-5.2.2| payload identifiers) are preserved. This documentation update does not change runtime behaviour or establish a new implementation release. Historical verification statements describe the tests recorded in the inherited baseline; this revision validates document compilation and references.

\section{Overview}
\label{sec:overview}

This document describes the \texttt{service\allowbreak{}Orchestration\allowbreak{}Plant\allowbreak{}Management} service, which provides management of the active engineered plant snapshot used by Dynamic\allowbreak{}Service\allowbreak{}Orchestration.
It is produced by the DynamicServiceOrchestration system \cite{dso-sysd} in service version 1.0.0.

The rest of this document is organized as follows.
In Section \ref{sec:operations}, we describe the abstract message operations provided by the service.
In Section \ref{sec:model}, we end the document by presenting the data types used by the mentioned operations.
The concrete interface is specified in the HTTP/TLS/JSON IDD \cite{dso-serviceOrchestrationPlantManagement-http-idd}.

\subsection{Significant Prior Art}

Local 2026-10-23 extension replacing the Plant\allowbreak{}Description\allowbreak{}Engine dependency on SimpleStore.

\subsection{How This Service Is Meant to Be Used}

PDE queries and atomically replaces the complete list of consumer/service scopes and ranked provider instance IDs. Lower nonnegative priority wins among otherwise eligible providers. LocalCloudManager may query, but only PDE or Sysop may replace, in addition to the common management policy.

\subsection{Important Delimitations and Dependencies}
\label{sec:delimitations}

HTTP only. Empty provider lists intentionally deny unrelated alternatives, including inter-cloud fallback. An empty rules list clears the active plant. DSO still enforces authorization, blacklist, locks, interface and QoS constraints. Existing consumer caches may delay visible changes. Apply database\_patch\_pde\_dso.sql before starting an upgraded DSO.

\newpage

\section{Service Interface}
\label{sec:operations}

This section describes the interfaces to the \texttt{service\allowbreak{}Orchestration\allowbreak{}Plant\allowbreak{}Management} service.
In particular, each subsection names an abstract operation, an input type and an output type, in that order.
The input type is named inside parentheses, while the output type is preceded by a colon.
Input and output types are only denoted when accepted or returned, respectively, by the interface in question.

All abstract data types named in this section are defined in Section \ref{sec:model}.

\begin{figure}[ht!]
  \centering
  \begin{tikzpicture}[font=\small]
    \node[draw=ArrowheadBlue, very thick, rectangle split, rectangle split parts=2, align=center,
          rectangle split part fill={ArrowheadBlue!12,white}, text width=9cm, inner sep=6pt] (b)
      {\guillemotleft{}Service\guillemotright{}\\\textbf{serviceOrchestrationPlantManagement}
        \nodepart{second} \texttt{query}\\ \texttt{replace}};
  \end{tikzpicture}
  \caption{SysML-style block of the \texttt{service\allowbreak{}Orchestration\allowbreak{}Plant\allowbreak{}Management} service and its operations.}
  \label{fig:SD}
\end{figure}

The following interface operations are available.

\osubsection{query}{}{\mref{OrchestrationPlantRules}}

Returns the current active-plant snapshot; allowed for authenticated PDE, LocalCloudManager or Sysop, also subject to the common management policy.

\osubsection{replace}{\mref{OrchestrationPlantRules}}{\mref{OrchestrationPlantRules}}

Atomically validates and replaces the complete snapshot. Only authenticated PDE or Sysop may replace; invalid input does not modify existing state. Empty rules clears the active plant; empty providers retains an unavailable managed scope.

\newpage

\section{Information Model}
\label{sec:model}

Here, all data objects that can be part of the \texttt{service\allowbreak{}Orchestration\allowbreak{}Plant\allowbreak{}Management} service are listed in alphabetic order.
Note that each subsection, which describes one type of object, begins with the \textit{struct} keyword, which is used to denote a collection of named fields, each with its own data type.
As a complement to the explicitly defined types in this section, there is also a list of implicit primitive types in Section \ref{sec:model:primitives}, which are used to represent things like hashes and identifiers.

\msubsection{struct}{OrchestrationPlantProvider}

Allowed service instance and its engineered priority.


\noindent\begin{tabularx}{\textwidth}{| >{\raggedright\arraybackslash}p{3.9cm} | >{\raggedright\arraybackslash\small}p{5.3cm} | X |} \hline
\rowcolor{gray!33} Field & Type & Description \\ \hline
service\allowbreak{}Instance\allowbreak{}Id & \pref{Identifier} & Service instance identifier \texttt{\textless{}System\textgreater{}\textbar{}\allowbreak{}\textless{}service\allowbreak{}Definition\textgreater{}\textbar{}\allowbreak{}\textless{}version\textgreater{}}. \\ \hline
priority & \pref{Number} & Nonnegative rank; zero is highest and lower values win among otherwise eligible providers. \\ \hline
\end{tabularx}



\msubsection{struct}{OrchestrationPlantRule}

Managed consumer/service scope, retained even without an available provider.


\noindent\begin{tabularx}{\textwidth}{| >{\raggedright\arraybackslash}p{3.9cm} | >{\raggedright\arraybackslash\small}p{5.3cm} | X |} \hline
\rowcolor{gray!33} Field & Type & Description \\ \hline
consumer & \pref{Name} & Name of the consumer system. \\ \hline
service\allowbreak{}Definition & \pref{Name} & Service definition name. \\ \hline
providers & \pref{List}$<$\allowbreak{}\mref[Orchestration\allowbreak{}Plant\allowbreak{}Provider]{OrchestrationPlantProvider}$>$ & Ranked allowed service instances; empty means managed but unavailable, with no unrelated or inter-cloud fallback. \\ \hline
\end{tabularx}



\msubsection{struct}{OrchestrationPlantRules}

Atomic snapshot owned by Plant\allowbreak{}Description\allowbreak{}Engine at DSO.


\noindent\begin{tabularx}{\textwidth}{| >{\raggedright\arraybackslash}p{3.9cm} | >{\raggedright\arraybackslash\small}p{5.3cm} | X |} \hline
\rowcolor{gray!33} Field & Type & Description \\ \hline
rules & \pref{List}$<$\allowbreak{}\mref{OrchestrationPlantRule}$>$ & Complete active-plant snapshot; an empty array clears all managed scopes. \\ \hline
\end{tabularx}



\subsection{Primitives}
\label{sec:model:primitives}

Types and structures mentioned throughout this document that are assumed to be available to implementations of this service.
The concrete interpretations of each of these types and structures must be provided by any IDD document claiming to implement this service.

\noindent\begin{tabularx}{\textwidth}{| p{4.6cm} | X |} \hline
\rowcolor{gray!33} Type & Description \\ \hline
\pdef{Identifier} & A string that identifies a record, e.g. a UUID or a composite identifier. \\ \hline
\pdef{List} & An \textit{array} of a known number of items, each having type A. \\ \hline
\pdef{Name} & A string identifier that is intended to be both human and machine-readable. \\ \hline
\pdef{Number} & An integer number. \\ \hline
\end{tabularx}

\newpage

\bibliographystyle{IEEEtran}
\bibliography{bibliography}

\newpage

\section{Revision History}
\subsection{Amendments}

\noindent\begin{tabularx}{\textwidth}{| p{1cm} | p{3cm} | p{2cm} | X | p{4cm} |} \hline
\rowcolor{gray!33} No. & Date & Version & Subject of Amendments & Author \\ \hline

1 & 2026-10-22 & 5.2.2 & Document created from the Arrowhead LaTeX templates for the 5.2.2 Java Spring implementation, including the local extensions of this workspace. & Jerker Delsing \\ \hline
2 & 2026-11-09 & 5.3.0 & New documentation edition; primary-source references and Go/Rust implementation scope; 5.2.2 retained. & Documentation maintenance \\ \hline
\end{tabularx}

\subsection{Quality Assurance}

\noindent\begin{tabularx}{\textwidth}{| p{1cm} | p{3cm} | p{2cm} | X |} \hline
\rowcolor{gray!33} No. & Date & Version & Approved by \\ \hline

1 & & \arrowversion  &  \\ \hline

\end{tabularx}

\end{document}
